% Extensión estimada: una página, cuatro párrafos (ficha 4.6).
% Fuentes: secciones 4.1 a 4.5. Sin números nuevos.
% Feedback del tutor (14/09/2026): el cierre se llama "Discusión"; la discusión que cruza clasificación y generación
% queda en el Capítulo 6.
\section{Discusión de los resultados de clasificación}
\label{sec:res_clasif_discusion}

Los resultados del capítulo permiten responder la primera pregunta de investigación
(Sección~\ref{sec:intro_objetivos}): en la prueba de \emph{HAM10000}, ningún \gls{vlm} sin una cabeza entrenada sobre
ese conjunto clasifica las lesiones tan bien como un clasificador supervisado. Los \glspl{vlm} generativos de propósito general quedan en el nivel del azar y
colapsan hacia una clase que depende del modelo y de la estrategia, y los controles muestran que ese comportamiento no
se explica por el desbalance de las clases ni por la modalidad de la imagen. DermLIP, en cambio, muestra que la
especialización en dermatología aporta una representación útil, pero solo se acerca a los supervisados con una cabeza
entrenada sobre \emph{HAM10000}. Las diferencias con los modelos sin ajuste son, además, significativas en McNemar,
salvo la de DenseNet frente a DermLIP en sensibilidad macro, y coinciden con lo que informa la literatura.

Entre los clasificadores supervisados, en cambio, la arquitectura pesa menos que otros factores. Con el mismo
procedimiento, Swin-Tiny, ConvNeXt-Tiny y ViT-B/16 no se distinguen entre sí, mientras que la semilla cambia la
sensibilidad de las clases minoritarias y del melanoma en más de una décima. En melanoma, la regla de decisión pesa más
que el modelo: el umbral fijado en validación lleva la detección de tres a cuatro de cada cinco casos sin cambiar el
\gls{auc}. Del mismo modo, la partición cambia el resultado, ya que una partición por imagen sobreestima la
sensibilidad macro en unas 0.07. Por ello, una comparación con otros trabajos tiene sentido solo cuando se conocen
su partición y su regla de decisión.

Estos resultados valen dentro de la fuente de entrenamiento. En \emph{HIBA}, Swin-Tiny pierde cerca del 40\,\% de su
sensibilidad macro, también en las imágenes dermatoscópicas, y el ruido de la imagen lo afecta más que la compresión
o la pérdida de resolución. Sus probabilidades, en cambio, están bien calibradas cuando se entrena sin pesos por clase.
En consecuencia, Swin-Tiny es la referencia del trabajo dentro de \emph{HAM10000}, pero no una herramienta validada
para otra población, y las limitaciones que eso implica, junto con las de evaluar las pruebas con una sola semilla, se
retoman en el Capítulo~\ref{cp:discusion}.

Estos resultados determinan también de dónde sale el diagnóstico del \emph{pipeline}. Como un \gls{vlm} generativo no
distingue las siete clases, el diagnóstico no puede salir del generador, y por eso lo fija Swin-Tiny, que se lo
entrega al generador a través de la cascada (Sección~\ref{sec:met_cascada}). Como el clasificador ocupa menos de
1~GB de memoria y clasifica una imagen en milisegundos, el costo del pipeline lo fija el generador. El
Capítulo~\ref{cp:res_gen} evalúa si, con ese diagnóstico, el reporte generado afirma la clase correcta y cuánto cuesta
generarlo, y la discusión que cruza ambos resultados queda para el Capítulo~\ref{cp:discusion}.
